% Part: many-valued-logic
% Chapter: three-valued-logics
% Section: introduction

\documentclass[../../../include/open-logic-section]{subfiles}

\begin{document}

\olfileid{mvl}{thr}{int}

\olsection{పరిచయం}

$\True$, $\False$లకు మరొక విలువ~$\Undef$ను మాత్రమే చేరిస్తే
మూడు-విలువల తర్కం వస్తుంది. అదనపు సత్యమూల్యం ఒక్కటే
అయినప్పటికీ, $\lnot$, $\land$, $\lor$, $\lif$లకు
సత్యమూల్య ప్రమేయాలను నిర్వచించగల విధానాలు చాలా ఉన్నాయి.
ఒక తర్కం ఆ సత్యమూల్య ప్రమేయాల ఏ కలయికనైనా వాడవచ్చు;
అలాగే $\True$ను మాత్రమే నిర్దేశిత విలువగా ఉంచాలా, లేదా
$\True$, $\Undef$ రెండిటినీ నిర్దేశిత విలువలుగా ఉంచాలా
అనే ఎంపిక కూడా ఉంటుంది.

ఇక్కడ ప్రసిద్ధ మూడు-విలువల తర్కాల్లో కొన్నింటిని,
వాటి ప్రేరణలను, కొన్ని లక్షణాలను పరిచయం చేస్తాం.

\end{document}
